\documentclass[11pt]{article}
\usepackage[margin=1in]{geometry}
\usepackage{amsmath,amssymb,amsthm}
\usepackage[hidelinks,pdfauthor={Bowen Zuo},pdftitle={Interlacing of excedance polynomials with a minimum cycle length}]{hyperref}
\newtheorem{theorem}{Theorem}
\newtheorem{lemma}[theorem]{Lemma}
\newcommand{\exc}{\operatorname{exc}}
\newcommand{\ii}{\mathrm{i}}
\allowdisplaybreaks[2]
\title{Interlacing of excedance polynomials with a minimum cycle length}
\author{Bowen Zuo\\[2pt]\small East China University of Science and Technology}
\date{November 2026}
\begin{document}
\maketitle
\begin{abstract}
For every integer cutoff \(k\ge3\) and upper size \(n\ge10^6(k-1)^2\), adjacent strict-excedance polynomials of derangements have strictly interlacing simple negative zeros after removal of their common simple zero. The same independent weights on cycle lengths two through \(k-1\) appear in both rows, with unit weights on longer cycles. For the five-weight cube at cutoff seven we lower this sufficient threshold from \(36{,}000{,}000\) to \(96\). We prove the unweighted result at every size for cutoffs four through six and give independent all-size proofs at cutoffs three and seven. The proof uses the nearest conjugate poles, an exact positive-label representation in the exterior, and exact finite certificates. We provide Lean~4 sources and a verification record.
\end{abstract}

\section{The statements}
For a permutation \(\sigma\) on \([j]\), put \(\exc(\sigma)=\#\{i:i<\sigma(i)\}\), and let \(m_2(\sigma)\) be its number of cycles of length two. Set
\[
 F_j(t,u)=\sum_{\substack{\sigma\in\mathfrak S_j\\ \sigma(i)\ne i\ \forall i}}
       t^{\exc(\sigma)}u^{m_2(\sigma)},\qquad
 P_j(t,u)=F_j(t,u)/t\quad(j\ge3).
\]
Taking \(u=0\) retains exactly the permutations with every cycle of length at least three. Alexandersson asks about real-rootedness and consecutive interlacing for every fixed minimum cycle length \(k\ge3\) \cite{alex}. The recurrence for these rows follows also by summing the finer enumeration of Rakotoson and Rakotondrajao \cite[Theorem~30]{rr}. Brenti proves real-rootedness on individual conjugacy classes and unimodality results for restrictions on cycle lengths \cite[Theorems~3.2--3.4]{brenti}. Ma \cite[Theorem~3.5]{ma} proves a Sturm property for a common weight on the cycle count.
For a cutoff \(k\), let \(E_{j,k}(t)\) denote the sum of \(t^{\exc(\sigma)}\) over permutations whose cycles all have length at least \(k\).
In contemporaneous work, ProgramComputer \cite{programcomputer,programcomputer3} proves the unweighted \(k=3\) case at every size and gives bounds for eventual unweighted interlacing. Their two-cycle theorem \cite{programcomputer2} covers upper sizes through \(3k-1\). Their later three-cycle theorem \cite{programcomputer4} covers upper sizes through \(4k\), including the first transition permitting four cycles. Their exact prefix and sharper uniform theorem \cite{programcomputer5,programcomputer6,programcomputer7} cover every admissible hard cutoff through upper size \(251\) and again from upper size \(10{,}001\).

Their aspect-at-least-twenty theorem \cite{programcomputer8} applies for \(k\ge7\) and lower size \(m\ge20k\). At \(k=7\) this starts at upper size \(141\), so together with their exact prefix it covers the hard cutoff seven at all sizes. Their subsequent proposed theorem \cite{programcomputer9} treats every lower size \(m\ge k\ge7\), with a linked computational review. We give an independent all-size proof for hard cutoff seven. For the shared profile of five independent weights, Theorem~\ref{seven} lowers the general threshold in Theorem~\ref{general} from \(36{,}000{,}000\) to \(96\).
\begin{theorem}\label{general}
Let \(k\ge3\) be an integer and let \(n\ge10^6(k-1)^2\) be the upper size. Choose the same weights \(w_2,\ldots,w_{k-1}\in[0,1]\) in both rows, with \(w_1=0\) and \(w_a=1\) for \(a\ge k\). Define the weighted row at size \(j\) by assigning each derangement the product of its cycle weights in its strict-excedance polynomial. The quotients of the rows at sizes \(n-1,n\) by \(t\) have respectively \(n-3,n-2\) simple negative roots, with exactly one lower-row root in every consecutive upper-row gap. Each original row has one further simple root at zero, and the full rows weakly interlace. In particular the conclusion holds for \(E_{n-1,k},E_{n,k}\), by taking all short weights zero.
\end{theorem}
The proof of Theorem~\ref{general} is in Section~\ref{sec:general}. For minimum cycle length three, the finite range below the analytic threshold can be completed.
\begin{theorem}\label{main}
For every \(n\ge4\), the roots of \(P_{n-1}(t,0)\) and \(P_n(t,0)\) are simple and negative. If the upper row has roots \(\beta_1<\cdots<\beta_{n-2}\) and the preceding row has roots \(\alpha_1<\cdots<\alpha_{n-3}\), then \(\beta_i<\alpha_i<\beta_{i+1}\) for all \(1\le i\le n-3\). Both full polynomials have one further simple root at zero and weakly interlace.
\end{theorem}
\begin{theorem}\label{family}
For \(n\ge512\), all the conclusions of Theorem~\ref{main} hold with any common mark \(u\in[0,1]\) in the two rows.
\end{theorem}
\begin{theorem}\label{four}
For minimum cycle length four the consecutive quotients \(E_{n-1,4}/t,E_{n,4}/t\) strictly interlace with simple negative roots at every upper index \(n\ge5\). The full rows weakly interlace. For \(n\ge512\) the same conclusions hold when derangements receive independent weights \(w_2,w_3\in[0,1]\) for their cycles of lengths two and three, respectively, and unit weight for every longer cycle.
\end{theorem}
\begin{theorem}\label{five}
At every upper size \(n\ge6\), the two quotients \(E_{n-1,5}/t,E_{n,5}/t\) have simple negative zeros in strict interlacing order. The full rows weakly interlace. For \(n\ge49\) the same conclusions hold throughout the closed cube of independent cycle weights \(w_2,w_3,w_4\), with \(w_1=0\) and unit weight on every cycle of length at least five.
\end{theorem}
Taking \(w_4=1\), or \(w_3=w_4=1\), also gives the weighted conclusions at cutoffs four and three from upper size 49.
\begin{theorem}\label{six}
At every upper size \(n\ge7\), the quotient rows \(E_{n-1,6}/t,E_{n,6}/t\) have respectively \(n-3,n-2\) simple negative roots, with exactly one lower-row root in each upper-row gap. Each full row has one further simple root at zero, and the full rows weakly interlace. The same conclusions hold for every upper size \(n\ge49\) throughout the closed cube of independent weights \(w_2,w_3,w_4,w_5\in[0,1]\), with \(w_1=0\) and \(w_a=1\) for every \(a\ge6\). The two rows use the same weights.
\end{theorem}
The first nonempty positive-size hard-cutoff row is \(E_{6,6}=tA_5\). In the weighted assertion, setting \(w_5=1\) gives the cutoff-five cube and also permits the preceding specializations.
\begin{theorem}\label{seven}
At every upper size \(n\ge8\), the quotient rows \(E_{n-1,7}/t,E_{n,7}/t\) have respectively \(n-3,n-2\) simple negative roots, with exactly one lower-row root in each upper-row gap. Each full row has one further simple root at zero, and the full rows weakly interlace. The same conclusions hold for every upper size \(n\ge96\) throughout the closed cube of independent weights \(w_2,\ldots,w_6\in[0,1]\), with \(w_1=0\) and \(w_a=1\) for every \(a\ge7\). The two rows use the same weights.
\end{theorem}
The initial nonempty hard-cutoff row is \(E_{7,7}=tA_6\). The weighted conclusion at cutoff six in Theorem~\ref{six} already applies from upper size 49.
We prove Theorem~\ref{family} by a uniform sign test at actual roots. Section~\ref{finite} completes Theorem~\ref{main} using finite integer inequalities. For \(j\ge3\), the coefficient of \(t\) in \(F_j\) is one since every cycle contributes at least one excedance and there is exactly one \(j\)-cycle having just one excedance. For example, the latter assertion follows from the single-cycle polynomial \(tA_{j-1}(t)\), with \(A_m\) the ordinary Eulerian polynomial. Inverting a fixed-point-free permutation sends each excedance to a deficiency and preserves its cycle type. Thus \(P_j\) is reciprocal of degree \(j-2\), monic, has constant term one, and is positive on \([0,\infty)\).

\section{Transport of a root sign}\label{transport}
Let \(P,Q\) be monic of degrees \(d,d+1\). Assume their coefficients vary continuously on an interval and both polynomials are positive on \([0,\infty)\). Suppose at every real root \(r\) of \(P\) we have \(Q(r)P'(r)<0\). If \(P\) splits at one parameter, it splits throughout: splitting is closed, since the ordered real roots of a convergent sequence of monic polynomials remain bounded and a subsequential limit gives a full factorization. It is open here because every real root is simple and a disjoint sign-change interval around each persists under sufficiently small coefficient changes. At any parameter the signs of \(Q\) alternate at the \(d\) ordered roots. The signs at \(-\infty\) and at zero give the two exterior roots, hence all \(d+1\) roots, with multiplicity one.

At \(u=1\) the ordinary derangement rows supply the splitting anchor. Define the auxiliary initial row \(P_2(t,1)=F_2(t,1)/t=1\). Then \(P_3(t,1)=1+t\), and for \(n\ge4\) their insertion recurrence is
\[
 P_n=(1+(n-2)t)P_{n-1}+t(1-t)P'_{n-1}+(n-1)tP_{n-2}.
\]
If \(P_{n-2}\) strictly interlaces \(P_{n-1}\), its sign at any root of \(P_{n-1}\) equals the sign of the derivative there. The two surviving terms thus both have sign opposite to that derivative, giving strict interlacing by induction.

For \(t=-e^{-x}\) write \(H_j=F_j(t,u)/(1-t)^j\), and let \(p=n-1\). Primes on \(H_p\) and \(P_p\) denote differentiation with respect to \(x\) and \(t\), respectively. At a predecessor root, \(-H'_p=t^2P'_p(t,u)/(1-t)^p\). The sign
\begin{equation}\label{goal}
 H_n(-H'_p)>0\qquad(H_p=0,\quad x\ge0)
\end{equation}
therefore implies the required polynomial sign on \([-1,0)\). If \(t<-1\), reciprocal symmetry gives \(P'_p(t)=-t^{p-4}P'_p(1/t)\) at the root and \(P_n(t)=t^{n-2}P_n(1/t)\). Their product changes by a positive factor. Thus \eqref{goal} and the anchor imply Theorem~\ref{family}.

\section{The normalized marked row}\label{sec:normalized}
The exponential generating function for all permutations by excedances is \((1-t)/(\exp((t-1)v)-t)\). By the labelled cycle formula, deleting fixed points multiplies it by \(e^{-v}\); marking two-cycles multiplies it by \(\exp((u-1)tv^2/2)\). Substitute \(v=z/(1-t)\), and set
\[
 L(z)=(1+e^{z-x})^{-1},\quad B(z)=e^{az+bz^2},\qquad
 a=\frac{-t}{1-t},\quad b=\frac{(u-1)t}{2(1-t)^2}.
\]
Then \(G(z)=\sum H_j z^j/j!=(1-t)B(z)L(z)\).
The bounds \(0\le a\le\min(1/2,e^{-x})\) and \(0\le b\le\min(1/8,e^{-x}/2)\) hold for \(x\ge0\). Since \((\partial_x+\partial_z)L=0\), direct differentiation gives
\begin{align}
 H_n+H'_p&=p\alpha H_{n-2}+p(p-1)\beta H_{n-3},\label{insert}\\
 \alpha&=ut/(1-t)^2,\qquad \beta=(1-u)t(1+t)/(2(1-t)^3).\nonumber
\end{align}
In particular \(|\alpha|,|\beta|\le e^{-x}\). Set
\[
 \zeta=x+\ii\pi=\rho e^{\ii\theta},\quad c=x/\rho,\quad s=\pi/\rho,\qquad
 Y_j=\frac{\rho^{j+1}H_j}{(1-t)j!},\qquad V_j=\rho^{j+1}[z^j]L(z).
\]
Write \(Z=\rho^{n+1}(H_n+H'_p)/((1-t)n!)\). Equation~\eqref{insert} reads
\[
 Z=(\alpha\rho^2Y_{n-2}+\beta\rho^3Y_{n-3})/n .
\]
The inequality \(|Z|<|Y_n|\) implies \eqref{goal} at any zero of \(H_p\).
\section{The three sign estimates}
Throughout this section \(n\ge512\) and \(0\le u\le1\).

\subsection{The first conjugate poles}
Assume \(0\le x\le16\), and put \(R=(x^2+4\pi^2)^{1/2}\). Then \(\rho<17,\ R<18,\ \rho/R<24/25\). The poles of \(L\) are \(x+\ii\ell\pi\), for odd signed integers \(\ell\), all with residue \(-1\). The contour \(|z|=R\) encloses precisely the two nearest ones. Its distance to each pole exceeds \(2/5\): the two smallest radial gaps are \(3\pi^2/(R+\rho)\) and \(5\pi^2/(R+\sqrt{x^2+9\pi^2})\). On the contour \(|1+e^{z-x}|>1/10\). Indeed, if \(v=\Re z-x\) has \(|v|\ge1\), use \(|1-e^v|\). Otherwise let \(w\in[-\pi,\pi]\) be the displacement of \(\Im z\) from the nearest odd multiple of \(\pi\). Then
\[
 |1+e^{z-x}|^2=(1-e^v)^2+4e^v\sin^2(w/2)
 \ \ge\ (v^2+w^2)/10.
\]
By \(aR\le xe^{-x}+\pi\) and \(bR^2\le x^2e^{-x}/2+\pi^2/2\), one has \(|B(z)|\le e^{17}\). At \(\zeta\), \(D=|B(\zeta)|\ge e^{-\pi^2/8}>1/4\). Moving the coefficient integral gives
\begin{equation}\label{poles}
 \left|Y_j-2\Re\{B(\zeta)e^{-\ii(j+1)\theta}\}\right|
 \le170\,e^{aR+bR^2}(\rho/R)^j <10^{-6}\qquad(j\ge509).
\end{equation}
For the last inequality split at \(x=8\). For \(x\le8\) one has \(\rho/R\le9/10\) by squaring, and \(170e^{17}(9/10)^{509}<10^{-6}\). For \(8\le x\le16\) the exponent is at most \((18+18^2/2)e^{-8}<1\), and \(170e(24/25)^{509}<10^{-6}\).

At \(Y_p=0\) let \(\phi=\arg B(\zeta)-n\theta\). Then \(|\cos\phi|\le2\cdot10^{-6}\) and \(|\sin\phi|>.99\). Since \(s\ge3/17\), the phase \(\phi-\theta\) in the next row and \eqref{poles} give \(|Y_n|>3Ds/2\). Both shorter rows have absolute value below \(3D\). For \(h=2,3\), \(\rho\le x+\pi\) and \(e^{-x}(x+\pi)^h\) decreases on the nonnegative half-line. Thus
\[
 |Z|\le 3D(\pi^2+\pi^3)/512 < (3D/2)(3/17).
\]

\subsection{The remaining central interval}
Assume \(16\le x\le3\sqrt{n-1}\). The partial fractions of the logarithmic derivative of the product for \(\cosh\) \cite{dlmf}, or its differentiated series, give
\[
 V_j=\sum_{\ell\ \mathrm{odd},\,\ell\in\mathbb Z}(c+\ii\ell s)^{-(j+1)}
 \quad (j\ge1),\qquad V_0=\rho/(1+e^{-x}).
\]
The differentiated sums converge absolutely and locally uniformly away from the poles. To control every order at once, write \(\ell=2i+1>0\). Each modulus of order at least two is at most \(\rho^2/(x^2+\pi^2\ell^2)\), which is at most
\[
 \rho^2\left(\frac1{x+2i}-\frac1{x+2i+2}\right),
\]
because \((x+\ell)^2-1\le2(x^2+\pi^2\ell^2)\). Telescoping and including conjugates gives \(|V_j|\le2\rho^2/x\le2(x+4)\) for \(j\ge1\); the constant term also satisfies the last bound. The nonnegative coefficients of \(B\) now give
\begin{align}
 |Y_j-V_j|&\le2(x+4)(e^{a\rho+b\rho^2}-1)
 \le4(x+4)^3e^{-x}<s/32.\label{bias}
\end{align}
Here \(a\rho+b\rho^2\le (x+4)^2e^{-x}<1/2\). For the last inequality use \(s\ge3/(x+4)\), the decrease of \(e^{-x}(x+4)^4\), and \(4\cdot20^4e^{-16}<3/32\).

For \(m\ge8,\ ms^2\ge1\) put \(U_m=V_{m-1}/2\). We claim
\begin{equation}\label{center}
 |U_m|\le1/8\quad\Longrightarrow\quad |U_{m+1}-cU_m|>s/2 .
\end{equation}
For positive odd \(\ell\), \(|c+\ii\ell s|^{-m}\le[1+(\ell^2-1)/8]^{-4}\), using \((1+h)^{m/8}\ge1+mh/8\). The sum of these bounds weighted by \(\ell\), over \(\ell\ge3\), is at most
\[
 3/16+4096(5^{-7}+5^{-6}/12)<1/3.
\]
Indeed the summand at 3 equals \(3/16\) in the bound and the remaining terms are at most \(4096\ell^{-7}\); compare the decreasing odd-lattice tail with an integral. The unweighted tail is below \(1/9\). Hence \(|\cos(m\theta)|\le1/8+1/9<1/4\). The \(\ell=1\) term of the centered difference has magnitude \(s|\sin(m\theta)|\). The remaining terms cost at most \(s/3\), since \(|(c+\ii\ell s)^{-1}-c|\le\ell s\) by squaring. This proves \eqref{center} under the stated assumptions \(m\ge8\) and \(ms^2\ge1\), for any \(x\ge0\).

In this interval \(ns^2\ge1\) since \(x^2\le9(n-1)\). At a zero of \(Y_p\), \eqref{bias} and \eqref{center} give \(|V_n-cV_p|>s\) and \(|Y_n|>15s/16\). Both shorter rows have absolute value below \(3(x+4)\), so
\[
 |Z|\le 6(x+4)^4e^{-x}/n<s/4.
\]
The last inequality uses the decrease of \(e^{-x}(x+4)^5\) and \(6\cdot 20^5e^{-16}/512<3/4\).

\subsection{The exterior}
Assume \(x\ge3\sqrt p\). Write \(r=e^{-x}\), \(w_h=h^p r^h\) for \(h\ge1\), and let \(m\) maximize these weights, with \(W=w_m\). If \(m>1\), comparison with \(m-1\) gives \(x\le p\log(m/(m-1))\le p/(m-1)\). Hence \(m\le2\sqrt p/5\) and \(m+1\le \sqrt p/3+2\le3\sqrt p/7\) for \(p\ge511\), also at \(m=1\). The product of the two neighbors divided by \(W^2\) equals \((1-m^{-2})^p\), with \(w_0=0\). The smaller neighbor is therefore below \(W/20\). Put \(\delta=1/200\). The ratio just past the right neighbor is at most \((1-(m+1)^{-2})^p\le e^{-49/9}<\delta\), by comparing successive ratios and using \(w_{m+1}\le W\). The left side has the same bound whenever that tail exists, and the ratios decrease outwards. Thus \(w_h\le W\delta^{|h-m|-1}\) at distances at least two. Summing yields
\[
 \sum_{h\ge1} w_h<4W,\qquad
 \sum_{h\ge1} h w_h\le m\sum w_h+\sum |h-m|w_h<6nW.
\]

Let \(T_j(h)=j![z^j]B(z)e^{hz}\). Expanding \(L\) geometrically near zero gives
\[
 H_j/(1-t)=\sum_{h\ge0}(-1)^h r^hT_j(h).
\]
For \(n-3\le j\le n\) and \(h\ge1\), expanding the tilt and bounding falling factorials by powers gives
\[
 1\le T_j(h)/h^j\le \exp(aj/h+bj^2/h^2)\le1+n^{-4}.
\]
Indeed \(x^2\ge9(n-1)\ge8n\) and \(32!\le8^{16}512^8\). Hence \(e^x\ge x^{32}/32!\ge n^8\) and \(r\le n^{-8}\). At \(h=0\) each monomial contributing to the coefficient uses at least \(\lceil j/2\rceil\) powers of \(r\); summing the inverse factorials of the linear and quadratic factors gives
\[
 0\le T_j(0)\le j!e^2 r^{\lceil j/2\rceil}\le n^{-4}r .
\]
For the last bound use \(\lceil j/2\rceil-1\ge j/3,\ x/3\ge2\log n\), and \(j!\le n^j\).
Put \(S_j=\sum_{h\ge1}(-1)^h r^hh^j\). As \(W\ge r\), the absolute error in replacing \(H_j/(1-t)\) by \(S_j\) is at most \(5n^{-4}W\) for \(n-3\le j\le p\) and at most \(7n^{-3}W\) for \(j=n\).

Let \(h_0\) be the label of the stronger immediate neighbor of \(m\), choosing the right one for \(m=1\). In \(S_n-h_0 S_p\) this neighbor disappears, while the mode contributes magnitude \(W\). Thus
\[
 |S_n-h_0 S_p|\ge W\left(1-2/20-2\sum_{d\ge2}(d+1)\delta^{d-1}\right)>4W/5 .
\]
At \(H_p=0\) the preceding error estimates and \(h_0\le n\) give \(|H_n|/(1-t)>3W/4\). Each of the two shorter rows has magnitude below \(5(1-t)W\); inserting these bounds in \eqref{insert} gives \(|H_n+H'_p|/(1-t)\le5(n+n^2)rW<W/4\). This proves \eqref{goal} also in the exterior and concludes the proof of Theorem~\ref{family}.

\section{An explicit threshold for independent weights}\label{sec:general}
Fix \(k\ge3\), put \(d=k-1\), and let \(w_1=0\), \(w_a\in[0,1]\) for
\(2\le a\le d\), and \(w_a=1\) for \(a>d\).
In this section \(F_j\) denotes the weighted strict-excedance row from
Theorem~\ref{general}. Use \(p=n-1\), \(r=e^{-x}\), and \(t=-r\), and retain
\(H_j,Y_j,V_j,\rho,c,s,Z\) from Sections~\ref{transport} and~\ref{sec:normalized}.
Throughout the section assume \(n\ge10^6d^2\). We prove \eqref{goal}
on the intervals cut at \(12d+100\) and \(3\sqrt p\).

\subsection{The normalized generating function and its insertion identity}
Define \(A_0(t)=1\) and
\[
 M_a=\frac{t A_{a-1}(t)}{(1-t)^a},\qquad
 b_a=\frac{(w_a-1)M_a}{a!},\qquad
 B(z)=\exp\!\left(\sum_{a=1}^{d} b_a z^a\right).
\]
The labelled cycle formula gives \(\sum H_j z^j/j!=(1-t)B(z)L(z)\).
For \(a\ge2\) the cycle polynomial is \(t A_{a-1}(t)\). For the linear
term, rewriting the unrestricted generating function in terms of \(L\)
produces an extra \(e^z\); deleting the fixed points multiplies this by
\(e^{-z/(1-t)}\), giving \(e^{b_1z}\).
The Eulerian recurrence gives \(M'_a=-M_{a+1}\), with differentiation
in \(x\). In particular \(b_1=-(\log(1-t))'\), and
\(b'_a+(a+1)b_{a+1}=(w_{a+1}-w_a)M_{a+1}/a!\), taking \(b_{d+1}=0\).
Since \((\partial_x+\partial_z)L=0\), coefficient extraction gives
\begin{align}
 C_a&=(w_{a+1}-w_a)M_{a+1}/a!,\nonumber\\
 H_n+H'_p&=\sum_{a=1}^{d}(p)_a C_a H_{p-a},
 &Z&=\frac1n\sum_{a=1}^{d} C_a\rho^{a+1}Y_{p-a},\label{general-insertion}
\end{align}
where \((p)_a=p!/(p-a)!\). The constant term of
\((\partial_x+\partial_z)\log((1-t)B)\) vanishes.
We use \(|Z|<|Y_n|\) to prove \eqref{goal} on the first two intervals.

Write \(v=V(\rho)\), \(V(R)=\sum_{a=1}^d |b_a|R^a\) and
\(J=\sum_{a=1}^d |C_a|\rho^{a+1}\). For every \(R\ge0\) the absolute
coefficient mass of the tilt satisfies
\[
 \sum_{i\ge0}|[z^i]B|R^i\le e^{V(R)} .
\]
Since the exponent of \(B\) has degree at most \(d\), each term in
\(V((1+1/d)\rho)\) is at most \((1+1/d)^d<3\) times the
corresponding term in \(V(\rho)\). Thus
\begin{equation}\label{quadratic-dilation}
 V((1+1/d)\rho)\le 3v .
\end{equation}

\subsection{Two bounds for the tilt}
For all \(x\ge0\) the pole formula in Section~4 gives, when \(a\ge2\),
\[
 \frac{\rho^a|M_a|}{(a-1)!}
 =|V_{a-1}|\le 2+\frac1{as^2}.
\]
The equality follows by the expansion of \(L\) at zero for \(x>0\)
and by continuity at \(x=0\). For the bound, write \(\ell=2b+1\ge3\);
then \(|c+\ii\ell s|^{-a}\le(2ab(b+1)s^2)^{-1}\) by Bernoulli's
inequality, and \(\sum_{b\ge1}1/(b(b+1))=1\). In particular this
bound is below 7 if \(x^2\le32a\), since \(\pi^2>9\). The bound
applies also at order \(a+1\) under the same condition. We will use
\begin{equation}\label{quadratic-tilt-cost}
 v\le11(1+\log d),\qquad J\le24d,\qquad
 x^2>32d\ \Longrightarrow\ v\le1+12d/x.
\end{equation}
For \(x\le8\) use the bound 7 for every nonlinear tilt order and
every insertion order. The linear tilt costs at most
\(\rho e^{-x}\le(x+\pi)e^{-x}<4\). This gives the first two estimates.
The last hypothesis is impossible when \(x\le8\).

For \(x\ge8\) put \(A=\min(d,\lfloor x^2/32\rfloor)\). Orders
\(a>A\) cost at most \(7\sum_{a>A}1/a\) in \(v\) and \(7(d-A)\)
in \(J\). For the low orders use
\(M_a=\sum_{h\ge1}(-r)^h h^{a-1}\), absolutely convergent for \(x>0\).
Their costs are bounded by
\[
 \sum_{h\ge1}\frac{e^{-hx}}h\sum_{a=1}^A\frac{(h\rho)^a}{a!},
 \qquad
 \rho\sum_{h\ge1} e^{-hx}\sum_{a=1}^A\frac{(h\rho)^a}{a!},
\]
respectively. Here \(x\le\rho\le9x/8\) and
\(\rho-x=\pi^2/(\rho+x)\le 8/x\). For \(h\le4A/x\), complete the
exponential and use \(h(\rho-x)\le32A/x^2\le1\). The head of
the first sum is at most \(3(1+\log d)\) and also at most
\(12A/x\). The head of the second is at most \(12\rho A/x\).
For \(h>4A/x\), the truncated exponential is at most
\(2^A e^{h\rho/2}\). Since \(\log 2<3/4\), inclusion of \(e^{-hx}\)
leaves a bound of \(e^{-hx/5}\).
The tails cost at most \(5/x\) and \(5\rho/x\), respectively.
These estimates give the first two bounds of \eqref{quadratic-tilt-cost},
using \(12\rho A/x+5\rho/x\le14A+6\). When \(x^2>32d\) one has
\(A=d\), and the bound \(1+12d/x\) follows from the low orders alone.

\subsection{A centered estimate in the compact interval}
Write \(z_\ell=(c+\ii\ell s)^{-1}\) for odd signed \(\ell\), and
\(D=|B(\zeta)|\), where \(\zeta=x+\ii\pi\). In particular
\(D\ge e^{-v}>0\). For \(\nu\ge8\),
\begin{equation}\label{quadratic-pole}
 \sum_{b\ge1}(2b+1)|z_{2b+1}|^\nu
 \le\frac{4}{(\nu s^2)^2}e^{-\nu s^2/2}.
\end{equation}
In fact, extract half the power from each summand. This part is at most
\((1+8s^2)^{-\nu/4}\le e^{-\nu s^2/2}\), by concavity of
\(u\mapsto\log(1+8u)\) on \([0,1]\) and \(\log9>2\).
For the other half apply
\((1+4b(b+1)s^2)^{\nu/8}\ge(\nu/2)b(b+1)s^2\) twice.
The resulting bound telescopes since
\((2b+1)/(b^2(b+1)^2)=b^{-2}-(b+1)^{-2}\).
Also \(|z_\ell-c|\le|\ell|s\), by the identity
\(z_\ell-c=s(s-\ii c\ell)/(c+\ii\ell s)\).

Put \(\Pi_j=2\Re(B(\zeta)z_1^{j+1})\).
We claim that
\begin{equation}\label{quadratic-centered}
 |Y_j-\Pi_j|< D/8\quad(n-d-1\le j\le n),\qquad
 |(Y_n-cY_p)-(\Pi_n-c\Pi_p)|<sD/8.
\end{equation}
To see this, write \(B_i=[z^i]B\).
The exact row is the convolution \(Y_j=\sum_{i=0}^{j}B_i\rho^i V_{j-i}\).
The principal term is the corresponding sum over all \(i\ge0\)
using the nearest pair, since \(z_1^{-1}=\zeta/\rho\).
Use the same head \(i\le T=\lfloor n/2\rfloor\) in both rows of the
centered difference. Each pole exponent in the ordinary head is
\(\nu=j-i+1\ge n/3\), since \(n\ge6d\); in the centered head it is \(\nu=n-i\ge n/3\),
with additional factor exactly \(z_\ell-c\).
In this interval \(\rho\le64d\), so \(\nu s^2\ge2\).
The ordinary head divided by \(D\), and the centered head divided by
\(sD\), are each at most \(2\exp(2v-ns^2/6)\), by
\eqref{quadratic-pole} and its conjugate. This is less than \(1/16\).
If \(x^2\le32d\) then \(\rho^2\le40d\), so
\(ns^2/6\ge (9\cdot10^6/240)d\), while \(2v\le22d\).
Otherwise write \(y=d/x\ge1/62\).
Then \(\rho^2\le2x^2\), giving \(ns^2/6\ge750000y^2\), and
\(2v\le2+24y\). The difference exceeds \(\log32\) at \(y=1/62\)
and is increasing thereafter.

For the tail, \(|V_j|\le\rho^2\) at all nonnegative orders, by the
pole bound above or the formula for \(V_0\). Every two-pole term
has modulus at most two before multiplication by \(B_i\rho^i\).
With \(R=(1+1/d)\rho\) we have
\[
 \sum_{i>T}|B_i|\rho^i\le e^{V(R)}(1+1/d)^{-(T+1)}.
\]
This includes the principal terms beyond the finite row convolution.
Therefore either tail divided by its respective scale has bound
\[
 \frac{4\rho^2}{s}\exp\!\left(V(R)+v-\frac{n}{2(d+1)}\right)<1/16.
\]
Indeed \(4\rho^2/s<350000 d^3\) and by \eqref{quadratic-dilation}
and \eqref{quadratic-tilt-cost} the positive exponent including the
prefactor and \(\log16\) is less than
\(23+3\log d+44(1+\log d)\le57d\).
But \(n/(2(d+1))\ge10^6d/3\).
This proves \eqref{quadratic-centered}.

At \(Y_p=0\), put \(S=B(\zeta) z_1^{p+1}\).
Then \(|\Re S|<D/16\), so \(|\Im S|>15D/16\). Since
\(\Pi_n-c\Pi_p=2s\Im S\), the centered estimate yields
\(|Y_n|>7sD/4\). The shorter rows in \eqref{general-insertion}
have magnitudes below \(17D/8\). From \(\rho\le64d\) and
\(J\le24d\) one gets \(J/s\le512d^2\). Thus \(|Z|<|Y_n|\).

\subsection{The intermediate interval}
On \(12d+100\le x\le3\sqrt p\), we have \(\rho\le4x/3\) and
\(ns^2\ge1\). In the positive-label sums above take \(A=d\) and
bound the whole truncated exponential by \(2^d e^{h\rho/2}\).
This yields
\[
 v\le2^{d+1}e^{-x/3},\qquad J\le\rho\,2^{d+1}e^{-x/3},\qquad
 \rho^3 2^{d+1}e^{-x/3}<1/32 .
\]
For the last bound use \(d\le(x-100)/12\), \(\log2<1\), and the
decrease of \(x^3e^{-x/4}\) on \([100,\infty)\).
The upper bound \(2(4/3)^3 100^3/(2^8 2^{25})\) is below \(1/32\).
Now \(|V_j|\le\rho^2/2\) by the telescoping bound of Section~4,
including order zero. The absolute tilt majorant gives
\[
 |Y_j-V_j|\le (\rho^2/2)(e^v-1)\le\rho^2v<s/32 .
\]
At \(Y_p=0\) the bare phase gap \eqref{center} gives
\(|V_n-cV_p|>s\), hence \(|Y_n|>15s/16\).
Since \(|Y_j|\le\rho^2\), the insertion is at most
\(\rho^2J/n<1/(32n)\le s/32\). This proves \eqref{goal}.

\subsection{The exterior from the exact positive labels}
For \(x\ge3\sqrt p\) we use \(p\ge\max(65536,4k^2)\), which
follows from \(n\ge10^6d^2\) and \(k=d+1\le2d\).
Write \(u=\sqrt p\). For a positive integer \(h\), choose a uniform
permutation of \(h\) slots and assign \(j\) labelled balls to these
slots independently and uniformly. If \(A_C\) is the number of
balls in a slot cycle \(C\), put \(\eta_j(h)=\mathbb{E}\prod_C w_{A_C}\),
with \(w_0=w_1=0\). The actual coefficient of \(t^h\) in
\(F_j(t)/(1-t)^j\) is \(h^j\eta_j(h)\). Indeed the cycle formula and
the positive-label expression for \(M_a\) write this row EGF as
\[
 \exp\!\left(\sum_{h\ge1}\frac{t^h}{h}
                   \sum_{a\ge2}w_a \frac{(hz)^a}{a!}\right).
\]
The coefficient of \(t^h\) is the cycle-index average over permutations
of the \(h\) slots. For each cycle of length \(b\), the factor
\((bz)^a/a!\) assigns each of the chosen \(a\) balls to one of its \(b\) slots;
\(j![z^j]\) completes the assignments of labelled balls.
At \(h=0\) the coefficient vanishes for \(j>0\).
In particular, with \(S_q=\sum_{h\ge1}(-r)^h h^q\),
\[
 H_j=\sum_{h\ge1}(-r)^h h^j\eta_j(h),\qquad
 -H'_p=\sum_{h\ge1}(-r)^h h^{p+1}\eta_p(h).
\]
Absolute local uniform convergence for \(x>0\) justifies differentiation.

Let \(L=\lfloor u/2\rfloor\). For \(h\le L\) and \(j\ge p\),
the mean \(\mu=j/h\) for each slot is at least \(2u\ge4k\).
When every slot has at least \(k\) balls, all its cycle weights are
one. For any individual binomial slot count \(X\),
Markov's inequality applied to \(2^{-X}\) bounds
\(\Pr(X\le\mu/4)\) by
\(\exp(-\mu/2+\mu\log(2)/4)\le e^{-\mu/4}\).
Thus
\begin{equation}\label{quadratic-slot-head}
 0\le1-\eta_j(h)\le h e^{-j/(4h)}
 \le(u/2)e^{-u/2}<1/(100n).
\end{equation}
We will also use \(e^{u/8}>200u^3\) for \(u\ge256\). It holds
at 256 since \(2^{32}>200\cdot256^3\), and
\(e^{u/8}/u^3\) is increasing there.

For \(q=p,n\), let \(m_q\) be a mode of \(h^q r^h\) over positive
\(h\) and let \(W_q=m_q^q r^{m_q}\).
Comparison with the preceding label gives \(m_q\le q/x+1\le u/3+2\).
For either exterior tail beyond the adjacent labels the weight at
distance \(b\ge2\) from the mode is at most
\(W_q 50^{1-b}\). To see this, the first such outward ratio is at
most \((1-(m_q+1)^{-2})^q\), by taking its ratio with the previous
one, which is at most one. On the left the same bound follows
with \(m_q-1\) in place of \(m_q+1\), whenever that tail exists.
Since \(m_q+1\le3u/7\), both bounds are at most
\(e^{-49/9}<1/50\); subsequent ratios decrease. Consequently
\[
 \sum_{h\ge1}h^q r^h<4W_q,\qquad
 \sum_{h>L} h^q r^h
 <2W_q 50^{3-u/6}\le2W_q e^{-u/8}<W_q/(100n).
\]
Here \(L-m_q\ge u/6-3\ge u/8\).
Combining \eqref{quadratic-slot-head} and the tail, using
\(0\le \eta_j(h)\le1\) everywhere and \(W_n\le nW_p\), gives
\begin{equation}\label{quadratic-exterior-errors}
 |H_p-S_p|\le W_p/(20n),\qquad
 |H_n-S_n|\le W_p/20,\qquad
 |-H'_p-S_n|\le W_p/20.
\end{equation}

Finally take \(h_0\) to be the stronger neighbor of \(m_p\),
with missing left neighbor given weight zero. The smaller neighbor
has weight below \(W_p/10\), since the product of the neighbor
weights at \(m_p>1\) is
\(W_p^2(1-m_p^{-2})^p\) and \(m_p\le 2u/5\).
Centering \(S_n-h_0S_p\) cancels the stronger neighbor and leaves
the mode with magnitude \(W_p\). The remaining immediate term is
bounded by \(2W_p/10\); the more distant terms cost at most
\(2W_p\sum_{b\ge2}(b+1)50^{1-b}\). Their total is strictly below
\(W_p/2\). It follows that
\[
 W_p/2<|S_n|+h_0|S_p|.
\]
At \(H_p=0\), \(h_0\le n\) and \eqref{quadratic-exterior-errors}
give \(|S_n|>9W_p/20\). Both true quantities in \eqref{goal}
have the sign of \(S_n\), proving the assertion in the exterior.

\subsection{The original rows}
For \(j\ge k\), the linear coefficient of the original row is one:
a single \(j\)-cycle has unit weight and its cycle polynomial is
\(tA_{j-1}(t)\), while two or more cycles give at least two excedances.
Inversion gives reciprocity of the quotient, which is monic with
constant term one and has no nonnegative real root.
For the chosen profile connect its short-cycle weights by a straight
segment to the all-one short-cycle vector, keeping the fixed-point
and longer-cycle weights fixed. The sign \eqref{goal} holds over the
whole segment at every predecessor zero in logarithmic coordinates.
Reciprocity and the transport argument of Section~\ref{transport}
give strict interlacing with simple negative quotient roots.
Restoring the common simple zero gives weak interlacing, and the
all-zero short-cycle vector yields the minimum-cycle-length
enumerators. This proves Theorem~\ref{general}.

\section{Finite rows at zero mark}\label{finite}
We describe the integer inequalities completing the proof for \(4\le n\le511\). Write \(P_n=P_n(t,0)\), \(y=t/(1+t)^2\), and \(P_n=(1+t)^{n-2}G_n(y)\). Put \(G_1=G_2=0,\ G_3=1\). The original insertion recurrence \cite[Theorem~30]{rr} specializes to
\[
 F_n=(n-1)tF_{n-1}+t(1-t)F'_{n-1}+
     \binom{n-1}{2}t(1+t)F_{n-3}\quad(u=0)
\]
with \(F_0=1,\ F_1=F_2=0\). Substitution after dividing by \(t\) gives for \(n\ge4\), with out-of-range coefficients set to zero,
\[
 [y^i]G_n=(i+1)[y^i]G_{n-1}
       +2(n-1-2i)[y^{i-1}]G_{n-1}
       +\binom{n-1}{2}[y^{i-1}]G_{n-3}.
\]
This follows by applying \((1+2(n-3)y)+y(1-4y)\,d/dy\) to the preceding gamma polynomial. The term of largest possible degree is positive: for even \(n\) its shift coefficient is \(2\), and for odd \(n\) the previous row already has that degree. In particular \(\deg G_n=\lfloor(n-2)/2\rfloor\) and \(G_n(0)=1\).

The supplementary certificate records disjoint negative dyadic intervals for every \(G_n\), in increasing order. To verify a row, evaluate its integer polynomial at both ends of each interval using homogeneous Horner evaluation, check opposite signs, and compare successive endpoints by integer arithmetic. There are as many intervals as the degree, so they exhaust the roots and each root is simple. For each even upper index, the upper row starts and ends the alternation of the two rows' intervals. For each odd upper index, the preceding row starts it and the upper row ends it. The recurrence, degree, bracket signs and ordering can be checked without repeating the search for intervals. The 508 pairs needed here use 64,770 brackets and 129,540 endpoint signs. The supplementary file also contains the longer independent check through size 1023. A second implementation formed the ordinary \(t\)-coefficients and inverted the triangular change of basis before evaluating the same brackets.

On \((-1,0)\), \(t\mapsto y\) is a strictly increasing bijection onto \((-\infty,0)\). Each gamma root gives one root on this interval and one reciprocal root on \((-\infty,-1)\). The term of highest gamma degree shows that at \(-1\) the even \(n\) row is nonzero and the odd \(n\) row has a simple root. For an even upper index the upper row lies closest to \(-1\) on either side and the predecessor has the central root; the roles reverse for odd upper indices. Thus the bracket alternation gives strict interlacing of the original quotients, including at the center. Restoring the common simple root at zero completes the proof of Theorem~\ref{main}.

\section{Minimum cycle length four}\label{sec:four}
We give the quantitative estimates for Theorem~\ref{four}. Use the notation of Section~\ref{sec:general} with \(d=3\). The two marks can vary throughout the closed square. The three coefficients of \(\log B(z)\) are
\[
 b_1=\frac r{1+r},\qquad
 b_2=\frac{(1-w_2)r}{2(1+r)^2},\qquad
 b_3=\frac{(1-w_3)r(1-r)}{6(1+r)^3}.
\]
They lie in \([0,1/2],[0,1/8],[0,1/24]\) and are also bounded by \(r,r/2,r/6\). In \eqref{general-insertion}, \(|C_1|,|C_2|,|C_3|\) satisfy the same latter bounds, since \(|1-4r+r^2|\le(1+r)^4\) on \([0,1]\). In particular the sign of \(w_3-w_2\) is unrestricted. In the three estimates below \(n\ge512\), \(p=n-1\), and the smallest coefficient index used is \(n-4\ge508\).

\paragraph{Compact interval.}
Write \(D=|B(\zeta)|\). Since \(\Re(\zeta^3)=x^3-3\pi^2x\ge-2\pi^3\), one has \(D\ge\exp(-\pi^2/8-\pi^3/12)>e^{-5}\).
The contour argument of \eqref{poles} bounds the absolute error in \(Y_j\) after retaining the nearest pole pair by \(170\exp(\sum b_i R^i)(\rho/R)^j\). The following upper bounds hold in the indicated order:
\[
\begin{array}{c|ccc}
x & R & \rho/R & \sum b_iR^i\\ \hline
{[0,4]}&8&7/10&34\\
{[4,8]}&11&9/10&6\\
{[8,16]}&18&24/25&1
\end{array}
\]
For the radii and ratios square both sides and use \(3<\pi<16/5\). For the first cost use the absolute bounds on \(b_i\); on the other intervals use \(r(R+R^2/2+R^3/6)\) with \(r\le e^{-4},e^{-8}\), respectively. At \(j\ge508\) each error is below \(10^{-6}<sD/16\), where \(s=\pi/\rho>3/17\). At \(Y_p=0\) the real part of \(B(\zeta)(\rho/\zeta)^n\) has absolute value below \(sD/32\). Rotating by \(\rho/\zeta\) thus gives \(|Y_n|>3sD/2\), while all three shorter rows have absolute value below \(3D\). For \(2\le h\le4\) one has \(r\rho^h\le\pi^h\), because the logarithmic derivative of \(e^{-x}\rho^h\) is at most \(-1+h/(2\pi)<0\). Hence
\[
 |Z|\le\frac{3D}{512}(\pi^2+\pi^3/2+\pi^4/6)
 < \frac9{34}D < \frac32sD .
\]

\paragraph{Middle interval.}
For \(16\le x\le3\sqrt p\), \(\rho\le33x/32\) and \(r\rho^6\le(33/2)^6e^{-16}<16/5\). Indeed \(x^6 e^{-x}\) decreases here. The tilt cost \(v=\sum b_i\rho^i\) and the drift sum \(J=\sum |C_i|\rho^{i+1}\) satisfy \(v\le(21/100)r\rho^3,\ J\le(21/100)r\rho^4\), using \(\rho\ge16\). The logistic coefficients satisfy \(|V_j|\le\rho^2/2\) for every \(j\), by the pole mass estimate (including \(V_0\)). Convolution gives \(|Y_j-V_j|\le(\rho^2/2)(e^v-1)\le\rho^2v<s/4\) as \(v\rho^3<27/40<\pi/4\). At \(Y_p=0\), \eqref{center} applies since \(ns^2\ge1\) and \(|V_p|<s/4\le1/4\); it yields \(|V_n-(x/\rho)V_p|>s\), so \(|Y_n|>s/2\). Every shorter row is bounded by \(3\rho^2/5\). Consequently \(n|Z|\le(3/5)(21/100)r\rho^6<1/2\le ns/2\).

\paragraph{Exterior.}
For \(x\ge3\sqrt p\) we have \(r n^8\le1\) as in the two-cycle case. In the expansion \(A_j=H_j/(1-t)=\sum_{h\ge0}(-r)^h T_j(h)\), one has \(T_j(h)=j![z^j]B(z)e^{hz}\). For \(h\ge1,\ j\le n\), the relative error \(|T_j(h)/h^j-1|\) is bounded by \(\exp(r(n+n^2/2+n^3/6))-1<n^{-4}\), using falling factorials bounded by the corresponding powers of \(j\). For the zero label take \(R_0=e^{x/3}\ge n^2\); Cauchy's estimate gives \(|T_j(0)|\le j!e^3 R_0^{-j}\le r n^{-4}\), because \(j+5\le2(j-3)\) for the indices here. Write \(W=\max_{h\ge1}h^p r^h\), \(S_j=\sum_{h\ge1}(-r)^h h^j\).
For \(p-3\le j\le n\) the moment estimates from the preceding exterior argument give \(|A_j-S_j|\le W/(8n)\): the positive-label sums are bounded by \(6nW\) for \(j\le p+1\), and \(r\le W\). At \(A_p=0\), subtracting the stronger neighbor index times \(S_p\) as before gives \(|A_n|>W/4\). For the three preceding indices \(|S_j|\le W\) by the alternating unimodal sum, hence \(|A_j|<2W\). The insertion in raw coordinates is bounded by \(6W r n^3<W/4\). This completes the derivative sign test on the square. The same path to the all-one derangement row and the same reflection argument now give strict interlacing of the quotients throughout that square.

\paragraph{Remaining original rows.}
At \(w_2=w_3=0\) the original row recurrence has short-cycle term \(\binom{n-1}{3}t(1+4t+t^2)E_{n-4,4}\), since the six cyclic orderings on four given letters have excedance polynomial \(t+4t^2+t^3\). In the gamma coordinates of Section~\ref{finite}, start with \(G_1=G_2=G_3=0,\ G_4=1+2y\). For \(n\ge5\),
\[
 G_n=(1+2(n-3)y)G_{n-1}+y(1-4y)G'_{n-1}
       +\binom{n-1}{3}y(1+2y)G_{n-4}.
\]
The constant term is one and the degree is \(\lfloor(n-2)/2\rfloor\) by the same top-coefficient argument. The supplementary dyadic certificate checks the 507 adjacent pairs between rows 4 and 511; counting the base row, these use 64,770 disjoint within-row brackets and 129,540 integer endpoint signs. One checker uses the gamma recurrence and the other reconstructs the ordinary rows before extracting their gamma coefficients. Degrees, interval order, reflection and the central root are handled exactly as in Section~\ref{finite}. This proves Theorem~\ref{four}.

\begingroup\raggedright
Theorem~\ref{four} uses \texttt{LongCycleBounds/FixedFourAllSizes.lean} for the original rows and \texttt{LongCycleBounds/FixedFourVerified.lean} for the weighted result. The finite witnesses and their Python checkers are in \texttt{certificates/k3/} and \texttt{certificates/k4/}. The local Lean build, kernel replay, and axiom audit records are indexed in \texttt{verification/README.md}.
\par\endgroup

\section{A fourth-degree tilt and cutoff five}\label{sec:five}
We prove Theorem~\ref{five}, retaining the notation of Section~\ref{sec:general}. In this section \(F_j\) denotes the independently weighted strict-excedance row. All comparisons use identical weights at the two sizes.
\paragraph{An adjustable two-pole estimate.}

For \(x\ge0\), put \(r=e^{-x},t=-r,\rho=|x+i\pi|,s=\pi/\rho,c=x/\rho\). Let \(L(z)=(1+r e^z)^{-1}\), \(B=\exp(\sum_{a=1}^d b_a z^a)\), \(B^\#(R)=\exp(\sum |b_a|R^a)\). Define \(V_j=\rho^{j+1}[z^j]L\), \(Y_j=\rho^{j+1}[z^j](BL)\). Put \(\zeta=x+i\pi\), \(z_1=\rho/\zeta\), \(D=|B(\zeta)|\). For every \(1<\lambda\le |x+3i\pi|/\rho\), \(R=\lambda\rho\), and every \(j\ge0\),
\[
 |Y_j-2\Re(B(\zeta)z_1^{j+1})|
 \le \max\!\left(\rho,\frac{\lambda\rho^2}{2\pi^2}\right)\lambda^{-j} B^\#(R).                 \tag{*}
\]
This bound applies to any real finite tilt, including either sign of its highest coefficient.

For \(m\ge1\), the logistic coefficients equal the absolutely convergent signed-odd pole sum, paired by conjugation. With \(z_h=\rho/(x+ih\pi)\), the nonprincipal remainder is at most \(2\sum_{l\ge1}|z_{2l+1}|^2\lambda^{1-m}\). Each squared modulus is at most \(\rho^2/[4\pi^2 l(l+1)]\). For \(m=0\), use the actual \(V_0=\rho/(1+r)\), with \(|V_0-2c|\le\rho\). Writing \(B(z)=\sum_i B_i z^i\), compare the finite convolution with the entire principal amplitude. For \(i<j\), the preceding remainder costs at most \(2\lambda\sum_{l\ge1}|z_{2l+1}|^2\) times \(\lambda^{-j}|B_i|R^i\). For \(i=j\) the corresponding constant is \(\rho\), and for \(i>j\) the missing principal amplitude costs at most \(2/\lambda\). The index ranges are disjoint and \(2/\lambda<\rho\), proving (*) by taking their maximum. This avoids paying the full tilt on a preselected circle of height \(2\pi\).

\paragraph{Returning to actual permutations.}

The cycle-profile exponential formula gives \(\sum H_j z^j/j!=(1-t)BL\), where \(H_j=F_j(t)/(1-t)^j\) and \(F_j\) is the original weighted row. Put \(M_a=t A_{a-1}(t)/(1-t)^a\). Then \(b_1=r/(1+r)\) and \(b_a=(w_a-1)M_a/a!\) for \(2\le a\le4\). The true derivative in \(x\), at fixed marks, satisfies
\[
H_n+H'_p=\sum_{a=1}^{4}(p)_a C_a H_{p-a},\qquad C_a=(w_{a+1}-w_a) M_{a+1}/a!.
\]
Consequently \(|Y_n|>n^{-1}\sum|C_a|\rho^{a+1}|Y_{p-a}|\) at \(H_p=0\) proves \(H_n(-H'_p)>0\). In unscaled raw coordinates \(A_j=H_j/(1-t)\) the insertion bound is instead \(\sum(p)_a|C_a A_{p-a}|\).

Write \(H=r/(1+r)^2\) and \(v=(1-r)/(1+r)\). Here \(A_m\) is the Eulerian polynomial and \((p)_a\) is the falling factorial. The four moments are \(-H,-Hv,-H(1-6H),-Hv(1-12H)\), for orders 2 through 5 respectively. All have magnitude at most \(r\). For the last term, \(v(1-12H)=3v^3-2v\). Its maximum on \([0,1]\) is one and its minimum is \(-4\sqrt2/9>-1\). The sign of \(b_4\) agrees with \(1-6H\) wherever the mark is below one; \(b_1,b_2,b_3\ge0\).



\paragraph{Compact, \(0\le x\le11\).} The rational checker divides this single real interval into 88 closed cells of length \(1/8\); it never samples marks or roots. In each cell it bounds the range of \(r,H,v\), of every coefficient, and of the real parts of \(\zeta^a\). From these it bounds \(B^\#(\lambda\rho)/D\), the geometrical factors in (*), and \(J=\sum |C_a|\rho^{a+1}/s\). Among the admissible rational choices \(\lambda\in\{21/20,22/20,\ldots,3\}\) it chooses the one giving the smallest proven cell bound. For all \(j\ge n-5\ge44\), its largest bound for the left side of (*) divided by \(sD\) is \(949/400000<3/1000\). Also \(J<48/5\). Set \(\epsilon=3/1000\). The imaginary part of \(S=B(\zeta)z_1^{p+1}\) at \(Y_p=0\) has magnitude at least \(D-|\Re S|\), and \(|\Re S|\le\epsilon sD/2\). Rotating once by \(c-is\) yields \(|Y_n|\ge(2-3\epsilon)sD\), whereas \(|Y_{p-a}|\le(2+\epsilon)D\). The strict inequality \((2+\epsilon)J/n<2-3\epsilon\) holds even using the global caps.
\paragraph{Middle, \(11\le x\le3\sqrt p\).} Write \(A=11\), \(\alpha=\sqrt{1+\pi^2/A^2}\), \(C=A^{-3}+(2A^2)^{-1}+(6A)^{-1}+1/24\). The tilt mass at \(\rho\) is at most \(C\alpha^4 x^4 e^{-x}\le v_0<0.017754\). For \(m\ge1\) the nonprincipal logistic poles have total absolute mass \(\le\rho^2/((m+1)\pi^2)\), by Bernoulli's inequality applied to \((1+4l(l+1)s^2)^{(m+1)/2}\) and summing \(1/[l(l+1)]\). For \(i\le j/2\), \(m=j-i\), this gives the bound \(2+2\rho^2/((n-3)\pi^2)\) on \(|V_m|\). The contribution of \(1\le i\le j/2\) to \(|Y_j-V_j|/s\) is at most
\[
 \frac{Ce^{-A}}{\pi(1-v_0)}\left(2\alpha^5 A^5+
 \frac{2\alpha^7 A^7}{(49-3)\pi^2}\right)<0.170,
\]
using monotonicity of \(x^h e^{-x}\) for \(h\le7\). For \(i>j/2\) bound \(|V_m|\le2\rho^2/x\) for all \(m\ge0\), including order zero. The absolute tilt at \(2\rho\) costs \(e^{16v_0}<2\). Since \(2\rho^3/(\pi x)<7n\), the remaining error divided by \(s\) is at most \(14n/2^{\lfloor(n-5)/2\rfloor}<0.000167\) by checking 49 and 50 and using monotonicity on each parity class. Thus \(|Y_j-V_j|<s/4\). At \(Y_p=0\) the already proved bare gap gives \(|V_n-cV_p|>s\) since \(ns^2\ge1\) and \(|V_p|<s/4\); hence \(|Y_n|>s/2\). For the insertion indices, \(|V_j|\le2+n/(n-4)\), so \(|Y_j|<4\). Finally \(J\le C\alpha^6A^6e^{-A}/\pi<0.740\), giving insertion below \(s/2\).
\paragraph{Exterior, \(x\ge3\sqrt p\).} For \(T_j(h)=j![z^j]B(z)e^{hz}\), \(A_j=\sum_{h\ge0}(-r)^hT_j(h)\). For all \(n\ge49\), \(r n^5<1/3\) and \(R_0=e^{x/4}>2n\). The inequalities at 49 are exact Taylor comparisons; their logarithmic derivatives give monotonicity. At \(h\ge1\), coefficient extraction gives \(|T_j(h)/h^j-1|\le\exp(r\sum_{a=1}^4 n^a/a!)-1<1/(60n)\), since the sum is below \(n^4/21\). At \(h=0\), using radius \(R_0\) gives \(|T_j(0)|\le6r n^4 2^{-(j-4)}\le r/(60n)\) for every needed \(j\ge n-5\); the tilt cost there is at most \(41/24\). At the zero of \(A_p\), put \(W=\max_{h\ge1}h^pr^h\), let \(h_0\) attain the maximum, and put \(S_j=\sum_{h\ge1}(-r)^h h^j\). The discrete peak bounds used in Section~\ref{sec:four}, with the size-$48$ peak estimate checked in the accompanying certificate, yield \(r\le W\), \(\sum_{h\ge1}h^jr^h\le4W\) for \(j\le p\), \(\sum h^nr^h\le4nW\), and \(W/2<|S_n|+(h_0+1)|S_p|\) with \(h_0+1\le n\). Thus \(|S_p|\le W/(12n)\) and \(|A_n-S_n|<W/12\), implying \(|A_n|>W/3\). The shorter alternating moments satisfy \(|S_j|\le W\) by unimodality, so \(|A_j|<2W\). The raw insertion is smaller than \(2Wr n^4/21<W/3\).

All three use closed regions and permit every point of the weight cube, including zero weights. The general homotopy readout passes through the derangement anchor, uses reciprocal symmetry to cover the second half of the negative axis, and gives the claimed strict interlacing of nonzero roots. Restoring the common simple zero gives the original weak interlacing.


\paragraph{Finite predecessor pairs.}\label{five-signed-remainder}
The two integer recurrences in the accompanying finite certificate checker use respectively the cycle containing a fixed letter, and deletion/insertion of the largest letter. They agree through 48; literal permutation enumeration separately agrees for sizes five through eight. Remove the simple factor \(t\) in each pair and begin with the upper, then the lower quotient. At each step replace the previous pair \((f,g)\) by \((g,-\operatorname{rem}(f,g))\), dividing by positive contents. Each remainder has the expected degree exactly one less than its predecessor and a positive leading coefficient. The terminal constant is positive. Backward induction using the opposite signs at zeros of \(g\) and the signs at infinity shows strict interlacing at every stage; it also proves simplicity. Positive original quotient coefficients exclude nonnegative zeros. The checker performs 1,032 exact remainder checks for the 43 target pairs and prints a digest of the normalized transcript.  This finite calculation concerns the hard cutoff alone.

\section{A fifth-degree tilt and cutoff six}\label{sec:six}

Use the notation of Sections~\ref{sec:normalized} and \ref{sec:general}, now with \(k=6\). We prove the insertion inequality \(|Z|<|Y_n|\) at a zero of \(H_p\) for every \(n=p+1\ge49\) and all four weights. Here
\[
 B(z)=\exp\!\left(\sum_{a=1}^5 b_a z^a\right),\qquad
 C_a=(w_{a+1}-w_a)M_{a+1}/a!,\qquad
 Z=\frac{1}{n}\sum_{a=1}^5 C_a\rho^{a+1}Y_{p-a}.
\]
Put \(u_0=r/(1+r)^2,\ v_0=(1-r)/(1+r)\), where \(r=e^{-x}\). The moments of orders two through six are
\[
 -u_0,\quad -u_0v_0,\quad -u_0(1-6u_0),\quad
 -u_0v_0(1-12u_0),\quad -u_0(1-30u_0+120u_0^2).
\]
For the last identity, the Eulerian polynomial is \(A_5(t)=1+26t+66t^2+26t^3+t^4\). The last quadratic takes values in \([-7/8,1]\) on \([0,1/4]\), as is seen from its vertex and endpoints. Also \(v_0(1-12u_0)=3v_0^3-2v_0\) has absolute value at most one on \(0\le v_0\le1\). Hence \(|M_a|\le u_0\) for \(2\le a\le6\), and
\[
 |b_a|\le r/a!,\qquad |C_a|\le r/a!\qquad(1\le a\le5).
\]
Write \(\mathcal V(R)=\sum_{a=1}^5 |b_a|R^a\) and \(J=\sum_{a=1}^5|C_a|\rho^{a+1}/s\).

\paragraph{Compact interval \(0\le x\le13\).}
The adjustable radius estimate \((*)\) of Section~\ref{sec:five} applies at \(R=\lambda\rho\) for \(1<\lambda\le|x+3\ii\pi|/\rho\). With \(D=|B(\zeta)|\ge e^{-\mathcal V(\rho)}\), its error divided by \(sD\) is at most
\[
 \max\!\left(\rho,\frac{\lambda\rho^2}{2\pi^2}\right)
 \frac{\rho}{\pi}\lambda^{-j}
       \exp\!\bigl(\mathcal V(\lambda\rho)+\mathcal V(\rho)\bigr).
\]
Partition \([0,13]\) into the 104 closed intervals \([h/8,(h+1)/8]\). On each interval bound \(r,u_0,v_0,\rho\) by their endpoint values, bound \(|1-6u_0|\) and \(|1-12u_0|\) by the greater endpoint magnitude, and use \(|M_6|\le u_0\). The attached rational certificate chooses a radius on each interval. It checks admissibility at the right endpoint, since \((x^2+9\pi^2)/(x^2+\pi^2)\) decreases, and bounds the exponential by rational Taylor enclosures. For \(j\ge n-6\ge43\) it gives an error less than \((7/50)sD\) and \(J<15\). The respective rational upper bounds in the arithmetic replay are \(0.1363569\) and \(14.2337303\).

Set \(\epsilon=7/50\) and \(S=B(\zeta)(\rho/\zeta)^{p+1}\). At \(Y_p=0\) we have \(|\Re S|\le\epsilon sD/2\) and \(|\Im S|\ge D-|\Re S|\). The pole term at \(n\) uses \((c-\ii s)S\). Its real part has magnitude at least \(sD-(c+s)|\Re S|\), so \(|Y_n|\ge (2-3\epsilon)sD\), using \(c+s\le2\). Each shorter term needed in the insertion formula has magnitude at most \((2+\epsilon)D\). The inequality \(15(2+\epsilon)<49(2-3\epsilon)\) proves the gap.

\paragraph{Middle interval \(13\le x\le3\sqrt p\).}
Set \(a_0=13\), \(\alpha=\sqrt{1+\pi^2/a_0^2}\), and
\[
 C=\sum_{a=1}^5 a_0^{a-5}/a! .
\]
Then \(\rho\le\alpha x\) and \(\mathcal V(\rho)\le Cr\rho^5\le v_\ast=13/1000\), because \(x^5e^{-x}\) decreases on this interval. Also \(\exp(\mathcal V(2\rho))\le e^{32v_\ast}<2\).
Write \(V_j=\rho^{j+1}[z^j]L(z)\), as before. For \(m\ge1\), the two nearest terms in the signed odd-pole sum cost at most two. For the remaining positive odd labels \(2h+1\), \(h\ge1\), Bernoulli's inequality gives
\[
 (1+4h(h+1)s^2)^{-(m+1)/2}
 \le \frac{1/h-1/(h+1)}{2(m+1)s^2}.
\]
Telescoping and counting the conjugates gives \(|V_m|\le2+1/((m+1)s^2)\). In the convolution for \(Y_j-V_j\), if the tilt index is \(1\le i\le j/2\) and \(j\ge n-6\), then \(2(j-i+1)\ge n-4\). Thus
\[
 |V_{j-i}|\le2+\frac{2\rho^2}{(n-4)\pi^2}.
\]
The nonconstant tilt mass is at most \(\mathcal V(\rho)/(1-\mathcal V(\rho))\). Consequently the part with \(i\le j/2\), divided by \(s\), is bounded above by
\[
 \frac{Ce^{-a_0}}{\pi(1-v_\ast)}
      \left(2\alpha^6 a_0^6+\frac{2\alpha^8 a_0^8}{45\pi^2}\right)
 <\frac{153}{1000}.
\]
Here \(n-4\ge45\) and \(x^m e^{-x}\) decreases for \(x\ge13\) and \(m\le8\). For the remaining convolution use the all-order bound \(|V_m|\le2\rho^2/x\) from the odd-pole sum. At \(m=0\), the explicit value \(V_0=\rho/(1+r)\) gives the same bound. Apply the tilt majorant at \(2\rho\). Since \(2\rho^3/(\pi x)<7n\), the rest divided by \(s\) has bound \(14n\,2^{-\lfloor(n-6)/2\rfloor}\le 0.0003272\), checked at 49 and 50 and decreasing in steps of two. Thus \(|Y_j-V_j|<s/4\) in every required order.

We have \(ns^2\ge1\). At \(Y_p=0\), \(|V_p|/2<s/8\le1/8\), so the bare phase inequality \eqref{center} gives \(|V_n-cV_p|>s\). The two error bounds then give \(|Y_n|>s/2\). The shorter rows satisfy \(|Y_j|<2+n/(n-5)+1/4<4\), by the same pole mass bound at order \(j\). Finally,
\[
 J\le Ce^{-a_0}\alpha^7 a_0^7/\pi<3/4 .
\]
Since \(4J/n<1/2\), the weighted insertion term is smaller than \(|Y_n|\).

\paragraph{Exterior \(x\ge3\sqrt p\).}
Write \(\mathcal H_j=H_j/(1+r)\), and set
\[
 T_j(h)=j![z^j]\bigl(B(z)e^{hz}\bigr),\quad
 \mathcal H_j=\sum_{h\ge0}(-r)^hT_j(h),\quad
 S_j=\sum_{h\ge1}(-r)^h h^j,\quad W=\max_{h\ge1}h^p r^h .
\]
The series for \(\mathcal H_j\) follows by expanding \(L\) geometrically in a neighborhood of \(z=0\), valid since \(x>0\). Let \(D(N)=\sum_{a=1}^5 N^a/a!\). For \(N\ge49\) we have \(D(N)\le N^5/100\) by dividing by \(N^5\) and evaluating the decreasing sum at \(49\). Also \(r n^6<14\). Indeed \((1+y^2)^6e^{-3y}\) decreases for \(y\ge69/10\), by logarithmic differentiation. A rational check at \(y=69/10<\sqrt{48}\) bounds its value by \(14\) using the first 64 terms of the series for \(e^{207/10}\). These bounds give \(rD(n)<1/(6n+1)\).

For \(h\ge1\), expanding the product and using \(j!/(j-i)!\le j^i\) gives
\[
 |T_j(h)-h^j|
 \le h^j(e^{\mathcal V(j/h)}-1)
 \le h^j/(6n) \qquad (1\le j\le n).
\]
At label zero, \(\mathcal V(2n)\le (32/100)rn^5<1/2\). For \(n-6\le j\le n\), the nonconstant coefficient bound at \(2n\) and \(j!\le n^j\) give
\[
 |T_j(0)|\le 2^{-j}(e^{\mathcal V(2n)}-1)
          \le rn^5 2^{-j}\le r/(6n).
\]
For the last comparison one can check \(12n^6\le 2^{n-6}\) at 49 and use \((50/49)^6<2\). In total
\begin{equation}\label{six-label-error}
 |\mathcal H_j-S_j|\le
     \frac{r+\sum_{h\ge1}h^j r^h}{6n}\qquad(n-6\le j\le n).
\end{equation}

We recall the exterior peak estimates with their constants. For \(q=p,n\), \(x\ge(29/10)\sqrt q\). The positive sequence \(h^q r^h\) is unimodal by its decreasing successive ratios. If \(m\) is a mode, then \(m-1\le q/x\) by comparison with the term at \(m-1\) when \(m\ge2\), so \(m\le\sqrt q/2\) for \(q\ge48\). The product of the two immediate neighbor ratios relative to the mode is at most \((1-1/m^2)^q\le e^{-4}<1/50\). The missing left neighbor at \(m=1\) contributes zero. The smaller neighbor is at most \(1/7\) of the peak. The ratio to the next farther term on the left is at most \(e^{-q/m^2}\). On the right the corresponding ratio is at most \(e^{-4}\) if \(m+1\le\sqrt q/2\), by contracting the preceding ratio by \((1-1/(m+1)^2)^q\). If \(m+1>\sqrt q/2\), it is at most \(e^{q/(m+1)-x}\le e^{-(9/10)\sqrt q}\). All farther ratios are therefore at most \(1/50\), giving total mass at most four times the peak.

With \(m_p\) a mode for \(p\), let \(h_0\) be its larger neighbor by mass. In \(S_n-h_0S_p\), the term at \(h_0\) vanishes and the term at \(m_p\) has magnitude \(W\). The remaining contributions have total magnitude at most
\[
 2W/7+2W\sum_{b\ge2}(b+1)50^{1-b}<W/2 .
\]
It follows that \(W/2<|S_n|+(m_p+1)|S_p|\). We have \(m_p+1\le n/8\) and \(m_n\le n/8\), because \(m_q\le\sqrt q/2\). At \(\mathcal H_p=0\), \eqref{six-label-error} gives \(|S_p|\le5W/(6n)\le W/n\). Also the peak in exponent \(n\) is at most \(m_nW\). Thus \(|\mathcal H_n-S_n|\le (1+4m_n)W/(6n)\le W/8\), and \(|\mathcal H_n|>W/4\).

For \(n-6\le j\le p\), an alternating sum of a positive unimodal sequence is bounded in magnitude by its peak, by splitting at the mode and grouping the pairs on either side. Thus \(|S_j|\le W\) and \eqref{six-label-error} implies \(|\mathcal H_j|<2W\). The unscaled insertion term, divided by \(1+r\), is at most
\[
 \sum_{a=1}^5(p)_a |C_a\mathcal H_{p-a}|
       \le 2WrD(n)<W/4 .
\]
This is the desired sign in the exterior.

\paragraph{Root order and the finite hard-cutoff range.}
The three bounds give \(H_n(-H'_p)>0\) at every zero with \(x\ge0\). The straight-line path to all unit weights on lengths at least two keeps the four short weights in the cube. Reciprocity and the derangement anchor in Section~\ref{transport} give the weighted claim. All short weights zero gives the hard cutoff in this range.

For upper sizes \(7\) through \(48\) use the \hyperref[five-signed-remainder]{signed-remainder test at the end of Section~\ref*{sec:five}}. The cutoff-six coefficient table can be generated either by the cycle containing the smallest label or by largest-label insertion, and these agree coefficient by coefficient. Literal permutation enumeration also checks sizes \(6\) through \(9\). Starting each quotient pair with the upper row, the negative-remainder chains, with division by positive contents, have \(1029\) degree-one drops and positive leading coefficients. Each chain ends in a positive constant. Backwards sign alternation gives the root order, and positive quotient coefficients place every root below zero. The dyadic endpoint certificate used in Lean gives the same conclusion directly by opposite signs in disjoint intervals and the exact degrees. This completes Theorem~\ref{six}.

\section{A sixth-degree tilt and cutoff seven}\label{sec:seven}

Retain \(r=e^{-x}\), \(p=n-1\), \(\rho=\sqrt{x^2+\pi^2}\), \(s=\pi/\rho\), and \(\zeta=x+\ii\pi\).
The main quantitative result of this section is the bound \(n\ge96\) with
the same five independent short-cycle weights in both rows.
For the cutoff-seven profile the normalized generating function and weighted
insertion formula from Section~\ref{sec:general} give
\begin{align*}
 B(z)&=\exp\!\left(\sum_{a=1}^{6}b_a z^a\right),
 & b_a&=(w_a-1)M_a/a!,\\
 Z&=\frac1n\sum_{a=1}^{6} C_a\rho^{a+1}Y_{n-a-1},
 & C_a&=(w_{a+1}-w_a)M_{a+1}/a!.
\end{align*}
Thus the shortest row in the insertion is \(n-7\). We prove
\(|Z|<|Y_n|\) at \(Y_p=0\) for all \(n\ge96\).

Write \(u_0=r/(1+r)^2\) and \(v_0=(1-r)/(1+r)\) as in Section~\ref{sec:six}.
The Eulerian recurrence gives \(A_6(t)=1+57t+302t^2+302t^3+57t^4+t^5\) and hence
\[
 M_7=-u_0v_0(1-60u_0+360u_0^2).
\]
This moment requires a new bound. At \(r=1/11\), \(|M_7|/r=258335/248832>1\).
In the variable \(v=v_0\in[0,1]\), let \(P(v)=17v/2-30v^3+45v^5/2\).
Then \(-M_7/u_0=P(v)\) and \(-M_7/r=(1+v)^2P(v)/4\).
For each of the 32 closed subintervals with endpoints \(h/32,(h+1)/32\),
apply the substitution \(v=(h+z)/32\). If the resulting power coefficients
are \(a_i\) in degree \(d\), the Bernstein coefficients are
\(\sum_{i=0}^{j}a_i\binom ji/\binom di\) for \(0\le j\le d\).
The Bernstein basis is nonnegative on \([0,1]\) and sums to one.
The maxima of the absolute coefficients are bounded by
\(152355739/83886080<11/6\) and \(1003440476713/962072674304<21/20\),
respectively. This proves \(|M_7|\le\min((11/6)u_0,(21/20)r)\), with the
endpoint \(r=0\) understood by continuity. All six tilt coefficients satisfy
\(|b_a|\le r/a!\) by the bounds through order six. For the drift,
\(|C_a|\le r/a!\) for \(a\le5\) and \(|C_6|\le(21/20)r/6!\).
Write \(\mathcal V(R)=\sum_{a=1}^{6}|b_a|R^a\) and
\(J=\sum_{a=1}^{6}|C_a|\rho^{a+1}/s\).

\paragraph{Compact interval \(0\le x\le16\).}
Use the adjustable two-pole bound \((*)\) from Section~\ref{sec:five},
at \(R=\lambda\rho\) with \(1<\lambda\le|x+3\ii\pi|/\rho\).
In each of the 128 closed eighth-grid intervals enclose \(r,u_0,v_0,\rho\)
by their endpoint bounds, using \(|M_6|\le u_0\) and both seventh-moment
bounds above. The rational certificate selects \(\lambda\) on the
grid of twentieths. It tests admissibility at the right endpoint,
where \((x^2+9\pi^2)/(x^2+\pi^2)\) is smallest. With \(D=|B(\zeta)|\),
the pole error divided by \(sD\) is at most
\[
 \epsilon=\max\!\left(\rho,\frac{\lambda\rho^2}{2\pi^2}\right)
 \frac{\rho}{\pi}\lambda^{-j}
 \exp\!\bigl(\mathcal V(\lambda\rho)+\mathcal V(\rho)\bigr).
\]
Use \(j\ge n-7\ge89\) and an upper bound for this expression on the cell.
The two-pole argument in Section~\ref{sec:six} yields
\(|Y_n|\ge(2-3\epsilon)sD\) at \(Y_p=0\) and
\(|Y_j|\le(2+\epsilon)D\) for the lagged rows.
Rational endpoint and exponential enclosures give a maximum of
\(\bigl(3\epsilon+(2+\epsilon)J/96\bigr)/2\) at most
\(246041/1250000<1\). The recorded upper bound allows every choice of
weights in the closed cube and proves the insertion inequality.

\paragraph{Middle interval \(16\le x\le3\sqrt p\).}
Let \(D_6(R)=\sum_{a=1}^6 R^a/a!\), \(K=D_6(16)/16^6=100117/47185920\)
and \(\alpha=51/50\). We have \(\rho\le\alpha x\),
\(\mathcal V(\rho)\le Kr\rho^6<1/200\), and
\(\exp(\mathcal V(2\rho))\le(1-64/200)^{-1}<2\).
Use \(157/50<\pi<22/7\), \(e^{-16}<1/8800000\), and the decrease
of \(x^m e^{-x}\) for \(x\ge16\) and \(m\le9\).
Split the tilt convolution at \(i=\lfloor j/2\rfloor\), for \(j\ge n-7\).
In the head, \(2(j-i+1)\ge n-5\), so the logistic bound of
Section~\ref{sec:six} gives \(2+2\rho^2/((n-5)\pi^2)\).
The total head error divided by \(s\) is therefore bounded by
\[
 \frac{K e^{-x}}{\pi(1-1/200)}
       \left(2\rho^7+\frac{2\rho^9}{(n-5)\pi^2}\right)
 <\frac1{16}.
\]
For the tail use \(|V_m|\le2\rho^2/x\) at all needed orders including zero.
Since \(2\rho^3/(\pi x)\le18\alpha^3p/\pi<7n\),
the doubled-radius bound makes the tail divided by \(s\) at most
\(14n\,2^{-\lfloor(n-7)/2\rfloor}<1/1000\), checked at \(n=96,97\)
and decreasing within each parity. Thus \(|Y_j-V_j|<s/4\).
We have \(ns^2\ge1\) since \(x^2\le9p<\pi^2p\).
At \(Y_p=0\) the bare gap \eqref{center} gives \(|V_n-cV_p|>s\),
where \(c=x/\rho\). The two errors give \(|Y_n|>s/2\), and all six
lagged rows satisfy \(|Y_j|<2+n/(n-6)+1/4<4\). Finally
\[
 J\le\frac{r}{\pi}\left(\sum_{a=1}^{5}\frac{\rho^{a+2}}{a!}
                        +\frac{21\,\rho^8}{20\cdot6!}\right)<2/5 .
\]
The last bound and the head bound follow by substituting
\(r\rho^m\le(\alpha\cdot16)^m/8800000\), \(n-5\ge91\), and the lower
bound on \(\pi\). The insertion costs at most \(4Js/n<s/2\).

\paragraph{Exterior \(x\ge3\sqrt p\).}
Use \(\mathcal H_j,T_j(h),S_j,W\) as in Section~\ref{sec:six}.
For \(n\ge96\),
\[
 rn^7<16,\qquad D_6(n)\le n^6/600,\qquad 6n^7<2^{n-7}.
\]
For the first inequality \(\sqrt{95}>487/50\), and the Taylor
sum through degree 96 at \(3(487/50)\) exceeds \(96^7/16\).
The function \(y^7e^{-3\sqrt{y-1}}\) decreases for \(y\ge96\)
since \(9y>14^2\). For the other two inequalities use the decrease of
\(D_6(n)/n^6\), and check the power-of-two bound at 96, then use
\((97/96)^7<2\).

For \(h\ge1\), \(j\le n\), coefficient extraction gives
\(|T_j(h)/h^j-1|\le e^{\mathcal V(n)}-1<1/(6n)\), since
\(\mathcal V(n)(6n+1)\le 16(6+1/n)/600<1\).
At \(h=0\) and \(n-7\le j\le n\) use the nonconstant tilt mass at \(2n\).
Here \(\mathcal V(2n)\le64rn^6/600<1/2\), and \(j!\le n^j\) gives
\[
 |T_j(0)|\le 2^{-j}\bigl(e^{\mathcal V(2n)}-1\bigr)
          \le 2^{-j}(128/600)rn^6<r/(6n).
\]
In the modal inequalities proved in Section~\ref{sec:six}, the
\(p\)-th absolute mass is at most \(4W\), the \(n\)-th at most
\(nW/2\), and
\(W/2<|S_n|+(m_p+1)|S_p|\) with \(m_p+1\le n/8\).
They apply to both indices because \(3\sqrt{n-1}\ge(29/10)\sqrt{n}\).
The two label bounds then give
\[
 |\mathcal H_j-S_j|<W/n\quad(n-7\le j\le p),\qquad
 |\mathcal H_n-S_n|<W/8 .
\]
At \(\mathcal H_p=0\) this implies \(|\mathcal H_n|>W/4\).
The shorter rows have \(|\mathcal H_j|<2W\) by unimodality of the
unsigned logistic weights. Their total contribution is at most
\(2W(21/20)rD_6(n)<W/4\). This proves the strict gap in the exterior.

\paragraph{The first 88 pairs for the independent hard-cutoff proof.}
The closed analytic intervals meet at \(x=16\) and \(x=3\sqrt p\).
Together they give \eqref{goal} throughout the five-dimensional cube
for every \(n\ge96\). The reciprocal and derangement arguments in
Section~\ref{transport} supply the original root order and simplicity.

For hard cutoff seven and \(8\le n\le95\), compute the two polynomial
rows independently by choosing the cycle through the smallest label
and by inserting the largest label. The first method uses the recurrence
\[
 E_{n,7}=\sum_{a=7}^{n}\binom{n-1}{a-1}tA_{a-1}(t)E_{n-a,7},
 \qquad E_{0,7}=1 .
\]
The second uses the insertion recurrence for the original rows from
\cite[Theorem~30]{rr}. Literal permutations check sizes 7, 8 and 9.
For each pair the certificate places all \(n-3\) lower roots in strictly
ordered negative dyadic intervals \([a_i,b_i]\), with opposite endpoint
signs. The upper row takes the same nonzero sign at \(a_i\) and \(b_i\),
alternating between consecutive intervals. It also has a certified
opposite sign to the left of the first interval, and its constant term
has sign opposite to the last bracket sign. Thus its \(n-2\) roots occur
strictly in the gaps and on both sides of the bracket list.
The exact degrees exhaust the roots of both rows and prove simplicity.
The data contain 4,268 lower-root brackets and 88 left-hand points.
This completes Theorem~\ref{seven}.

\paragraph{Formal verification.}
The Lean development uses the original weighted permutation rows and the minimum-cycle enumerators. The accompanying statement map gives the entry declarations for Theorems~\ref{general}--\ref{seven}. The verification record specifies the build and kernel replay status for each group of declarations.
\paragraph{Acknowledgment.}
ChatGPT assisted with the analytic derivations, the implementation and independent replay of the formal proofs, and preparation of the manuscript.
\begingroup
\small
\begin{thebibliography}{99}
\setlength{\itemsep}{0pt}
\bibitem{alex}
P.~Alexandersson, \emph{Real-rooted polynomials: open problems},
\href{https://www.symmetricfunctions.com/realRootedOpenProblems.htm\#permutations-with-long-cycles}{Permutations with long cycles},
The Symmetric Functions Catalog, 2026.
\bibitem{programcomputer}
ProgramComputer, \emph{Consecutive excedance zeros of permutations with a minimum cycle length}, November 2026,
\href{https://github.com/ProgramComputer/open-math-research/blob/4da48b72ed8f92f96c2faf60e2d63d253b3b507a/topics/long-cycle-excedances/paper.tex}{Theorems~1 and~2}.
\bibitem{programcomputer3}
ProgramComputer, \emph{A numerical cutoff for every minimum cycle length},
November 2026,
\href{https://github.com/ProgramComputer/open-math-research/blob/4da48b72ed8f92f96c2faf60e2d63d253b3b507a/topics/long-cycle-excedances/effective-interlacing.tex}{uniform unweighted bound}.
\bibitem{programcomputer2}
ProgramComputer, \emph{Consecutive interlacing through the two-cycle range}, November 2026,
\href{https://github.com/ProgramComputer/open-math-research/blob/4da48b72ed8f92f96c2faf60e2d63d253b3b507a/topics/long-cycle-excedances/two-cycle-insertion.tex}{two-cycle insertion theorem}.
\bibitem{programcomputer4}
ProgramComputer, \emph{Strict interlacing through the first transition to four cycles}, November 2026,
\href{https://github.com/ProgramComputer/open-math-research/blob/e554fb387edbf68f5e2094773016f81b9d618498/topics/long-cycle-excedances/three-cycle-complete.md}{three-cycle theorem}.
\bibitem{programcomputer5}
ProgramComputer, \emph{Exact long-cycle excedance interlacing through size 250}, November 2026,
\href{https://github.com/ProgramComputer/open-math-research/blob/e554fb387edbf68f5e2094773016f81b9d618498/topics/long-cycle-excedances/finite-prefix/README.md}{finite prefix}.
\bibitem{programcomputer6}
ProgramComputer, \emph{Consecutive excedance interlacing from size 10,000 at bounded ratios}, November 2026,
\href{https://github.com/ProgramComputer/open-math-research/blob/e554fb387edbf68f5e2094773016f81b9d618498/topics/long-cycle-excedances/paired-residues/README.md}{paired-residue estimate}.
\bibitem{programcomputer7}
ProgramComputer, \emph{A common-circle comparison without a height loss}, November 2026,
\href{https://github.com/ProgramComputer/open-math-research/blob/e554fb387edbf68f5e2094773016f81b9d618498/topics/long-cycle-excedances/fourier-circle/README.md}{Fourier estimate}.
\bibitem{programcomputer8}
ProgramComputer, \emph{Consecutive excedance interlacing at aspect at least twenty}, November 2026,
\href{https://github.com/ProgramComputer/open-math-research/blob/6db8b1ec8c75269898e64b8676fe12983f7f008c/topics/long-cycle-excedances/aspect-20/README.md}{research theorem associated with PR~\#214}.
\bibitem{programcomputer9}
ProgramComputer, \emph{A proposed proof for all minimum cycle lengths at least seven}, November 2026,
\href{https://github.com/ProgramComputer/open-math-research/blob/a9f329264e850ab4142880e5e68a0c2496cf4a34/topics/long-cycle-excedances/all-large-cutoffs.tex}{research manuscript associated with PR~\#218}.
\bibitem{brenti}
F.~Brenti, \emph{Permutation enumeration, symmetric functions, and
unimodality}, Pacific J. Math. \textbf{157} (1993), 1--28,
\href{https://msp.org/pjm/1993/157-1/pjm-v157-n1-p01-s.pdf}{Theorems~3.2--3.4 and Problem~5.5}.
\bibitem{ma}
S.-M.~Ma, \emph{Counting permutations by numbers of excedances, fixed points
and cycles}, Bull. Aust. Math. Soc. \textbf{85} (2012), 415--421,
\href{https://doi.org/10.1017/S0004972711002760}{Theorem~3.5}.
\bibitem{rr}
I.-N.~M.~Rakotoson and F.~Rakotondrajao,
\emph{Associated Stirling Eulerian numbers and their \(q\)-analogues},
J. Integer Seq. \textbf{28} (2025),
\href{https://cs.uwaterloo.ca/journals/JIS/VOL28/Rakotoson/rakoto12.html}{Article 25.5.7},
Theorem~30.
\bibitem{dlmf}
NIST Digital Library of Mathematical Functions,
\href{https://dlmf.nist.gov/4.36.E2}{Eq.~4.36.2}.
\end{thebibliography}
\endgroup
\end{document}
